\documentclass[11pt]{article}
\usepackage[T1]{fontenc}
\usepackage{lmodern,amsmath,amssymb,amsthm,mathtools,booktabs,array}
\usepackage[margin=1in]{geometry}
\usepackage{microtype,xcolor,fancyhdr,hyperref}
\definecolor{navy}{HTML}{173653}
\hypersetup{colorlinks=true,linkcolor=navy,urlcolor=navy,citecolor=navy,pdftitle={Beyond Dirichlet GRH: a sharp PSL2 transfer}}
\pagestyle{fancy}\fancyhf{}\lhead{\small\textsc{Research note}}\rhead{\small Beyond Dirichlet GRH}\cfoot{\thepage}\setlength{\headheight}{14pt}
\newtheorem{theorem}{Theorem}[section]
\newtheorem{lemma}[theorem]{Lemma}
\newtheorem{corollary}[theorem]{Corollary}
\newtheorem{proposition}[theorem]{Proposition}
\theoremstyle{remark}\newtheorem{remark}[theorem]{Remark}
\DeclareMathOperator{\PSL}{PSL}\DeclareMathOperator{\SL}{SL}\DeclareMathOperator{\PGL}{PGL}\DeclareMathOperator{\Ind}{Ind}\DeclareMathOperator{\ord}{ord}\DeclareMathOperator{\Gal}{Gal}\DeclareMathOperator{\Tr}{Tr}\DeclareMathOperator{\Rea}{Re}\DeclareMathOperator{\FP}{FP}
\newcommand{\Q}{\mathbb Q}\newcommand{\C}{\mathbb C}\newcommand{\F}{\mathbb F}\newcommand{\one}{\mathbf 1}\newcommand{\St}{\mathrm{St}}\newcommand{\Gg}{\Gamma}\newcommand{\PT}{\mathcal P}\newcommand{\eps}{\varepsilon}
\newcommand{\repo}{\texttt{grh-lean}}
\begin{document}
\begin{center}
{\small\color{navy}\textsc{Mathematical research \enspace|\enspace 25 November 2026}}\\[10pt]
{\LARGE\bfseries Beyond Dirichlet GRH}\\[4pt]
{\Large Sharp nonabelian transfer theorems for $\PSL_2(2^r)$}\\[8pt]
{\small Analysis of the Alp\"oge--Sernau Lean proof architecture}
\end{center}
\begin{abstract}
The Lean development \repo\ establishes the Riemann hypothesis for Dirichlet $L$-functions. We extract a different target for its one-sided prime-sum method: Dedekind-zeta RH for every Galois extension $M/\Q$ with group $\PSL_2(q)$, $q=2^r\ge4$. We construct a nonnegative character $W$ of dimension $q^2$ containing each irreducible exactly once, prove its dimension is minimal among such nonnegative detecting characters, and show that a single zeta function of degree $q(q+1)$, or two of degrees $q+1$ and $q(q-1)$, suffice for RH of $\zeta_M$. Both degree thresholds are sharp for detection by subfield permutation characters. No Artin holomorphy conjecture is assumed. We reduce the outstanding RH to a one-sided inequality counting rational primes with zero, one, two, or $q+1$ roots of a degree-$q+1$ polynomial modulo $p$. The inequality has the same smoothed kernel and exponent as the Lean proof, but its arithmetic establishment remains open. Exact finite group checks for $q=4,8,16$ accompany the general proofs; subgroup searches for $q=4,8$ are exhaustive.
\end{abstract}
\noindent\textbf{Status.} The transfer theorems, optimality statements and analytic equivalences are proved below using classical results explicitly identified. This is not a proof of nonabelian RH or of a further named major conjecture. It gives a more economical, precisely specified research target and rules out certain cheaper strategies. The results are independently derived here; worldwide novelty is not asserted.
\tableofcontents

\section{What the Lean proof supplies, and what it does not}
We take the requested RH and Dirichlet GRH as available. The upstream repository \cite{lean} at commit \texttt{72fead102f663d5c7a407964a13dc0cf9d977207} states RH for the Riemann zeta function and GRH for Dirichlet $L$-functions, with the usual qualifications about trivial zeros and poles. It does not state RH for general Dedekind zeta functions, Artin or automorphic $L$-functions. The published repository's Lean checks are evidence for its own statements, not for the results in this report. We did not build or kernel-check the upstream development locally.

The readable edition supplies a particularly useful architecture. The following references are paths relative to \texttt{grh/src/}; exact blob identifiers are in \texttt{source\_manifest.tsv}.
\begin{center}\small
\begin{tabular}{@{}>{\raggedright\arraybackslash}p{.42\linewidth}>{\raggedright\arraybackslash}p{.51\linewidth}@{}}
\toprule
Source & Role in the argument\\\midrule
\texttt{Statements.lean} & Endpoint and scope of GRH.\\
\texttt{New/Reduction/MainInequality.lean}, line 87 & \texttt{star\_inequality}: short forms, enough independent forms, and a lower bound on integral forms.\\
\texttt{New/Reduction/PrimeSums.lean}, lines 89--177 & Convert exact local divisibility and congruence-class counts into weighted sums.\\
\texttt{New/Reduction/OneSidedBound.lean}, lines 33, 200, 239 & Cancel main terms, choose parameters, and derive an eventual one-sided error bound.\\
\texttt{New/Reduction/Zeros/Criterion.lean}, line 37 & Exclude zeros using one-sided bounds, an auxiliary quadratic twist, and Landau's argument.\\
\texttt{New/Reduction/Zeros/PrimeError.lean}, lines 65, 414 & Smoothed remainder and its Mellin-transform poles.\\\bottomrule
\end{tabular}
\end{center}
The reduction takes a remainder built from weights of the form
\[
 \Lambda(n)(1-\chi_K(n))\one_{n\equiv1\pmod {q_0}},\qquad
 k_x(n)=(1/n-1/x)\one_{n\le x},
\]
subtracts the main logarithm and a constant, and bounds the resulting error from above. In the comparison of covolumes, one first obtains upper bounds from short forms and a lower bound from integrality; a sufficiently large independent family makes the main contribution cancel. In \texttt{star\_inequality} the bound has the shape $m\,\mathrm{rem}\le fa\log^5 b$. Choosing scales, in particular $m$ of order $\sqrt x$, leads to the exponent $-1/2+\eps$. The remaining step is a one-sided analytic zero-exclusion criterion. The cyclotomic and quadratic parts of this construction encode abelian Frobenius data. Replacing the weights is not merely changing a parameter in that proof.

Our approach keeps the analytic step, but derives a different finite-group input. We choose a positive detector so that its associated $L$-function has no unknown poles and detects every nontrivial representation. This removes the need for an auxiliary twist to avoid cancellation of zero orders at the analytic stage.

\paragraph{Relation to previous reports.} The supplied \texttt{ProveIt/docs/incoming} folder \cite{incoming} contains multiple GRH-transfer archives. We inspected the central-directory listings of eleven archives and read the principal TeX/README summaries of four smaller archives (listed in the manifest). These treat, among other topics, cubic and $A_5$ detectors, class groups, CM equations, arithmetic progressions, sparse sampling, and the impossibility of recovering genuinely nonabelian Frobenius data using fixed congruence classes. In particular, an earlier $A_5$ detector used weights $(21,1,0,1,1)$ with its own choice of constraints. Here we impose the distinct criterion that a genuine character contain every nontrivial irreducible, be pointwise nonnegative, and have smallest dimension. Our construction works uniformly for even-field $\PSL_2$ and gives explicit optimal subfield degrees. Two larger ZIPs were not fully accessible via the connector; this is not an exhaustive comparison with every earlier report.

\section{A small nonnegative detector for a large simple group}
\subsection{Notation and the four types of elements}
Let $q=2^r\ge4$ and $G=\PSL_2(q)=\SL_2(q)=\PGL_2(q)$ (canonical identifications in even characteristic). Its order is $g=q(q^2-1)$. Let $G$ act on $X=\mathbb P^1(\F_q)$, of size $q+1$. Denote by $B$ a point stabilizer, by $U\lhd B$ the additive translation group, by $T$ the split torus of order $q-1$, and by $T'$ a nonsplit torus of order $q+1$. Their normalizers $N,N'$ are dihedral of orders $2(q-1)$ and $2(q+1)$. The stabilizer $B$ has order $q(q-1)$ and $B/U\simeq\F_q^\times$.

Write $f(h)=|X^h|$. The conjugacy types in the following table follow from the characteristic polynomial of a determinant-one $2\times2$ matrix. The nontrivial semisimple classes are indexed by pairs of inverse eigenvalues, in $\F_q^\times$ or the norm-one subgroup of $\F_{q^2}^\times$. Centralizers of those matrices are the corresponding tori; the centralizer of a nonidentity unipotent is $U$. All nonidentity unipotents are conjugate (squaring is bijective on $\F_q^\times$). In particular, the number of conjugacy classes is $q+1$.
\begin{center}\small
\begin{tabular}{@{}lrrrr@{}}\toprule
Type & $f$ & Centralizer size & Classes & Elements\\\midrule
Identity & $q+1$ & $g$ & $1$ & $1$\\
Unipotent $u\ne1$ & $1$ & $q$ & $1$ & $q^2-1$\\
Split semisimple & $2$ & $q-1$ & $(q-2)/2$ & $q(q+1)(q-2)/2$\\
Nonsplit semisimple & $0$ & $q+1$ & $q/2$ & $q^2(q-1)/2$\\\bottomrule
\end{tabular}
\end{center}

\subsection{Explicit induced characters}
For $H\le G$, put $\pi_H=\Ind_H^G\one$. The projective-line character is $\pi_B=f=\one+\St$, with $\St$ irreducible of degree $q$: the action is doubly transitive, so $\langle f,f\rangle=2$ and $\langle f,\one\rangle=1$. Let $\lambda$ be any nontrivial character of the additive group $U$; its values are $\pm1$. Define
\[
 \Gg=\Ind_U^G\lambda,\qquad W=\one+\Gg.
\]
Define $P$ as the sum of $\Ind_B^G\theta$ for one representative of each inverse pair of nontrivial characters $\theta$ of $B/U$, inflated to $B$. Define $D=\Ind_{N'}^G\eta$, where $\eta$ is the character equal to $1$ on $T'$ and $-1$ on $N'\setminus T'$. These are genuine characters, not virtual differences. We use capital $P,D$ for characters, not primes or discriminants.

\begin{lemma}[Character identities]\label{lem:identities}
The following identities hold, with $\mathrm{reg}$ the regular character of $G$:
\begin{align}
 \Gg&=\St+P+D, & W&=\one+\St+P+D,\label{eq:gamma}\\
 \pi_T&=\one+2\St+P+D, & \pi_{T'}&=\one+P+D,\label{eq:perms}\\
 \mathrm{reg}&=\one+q\St+(q+1)P+(q-1)D.\label{eq:reg}
\end{align}
In particular $W=\pi_B+\pi_{T'}-\one$ and $\pi_T=W+\St$. Its values are $q^2,0,1,1$ on the four types above.
\end{lemma}
\begin{proof}
Use the induction formula
\[
 (\Ind_H^G\alpha)(h)=\frac{|C_G(h)|}{|H|}
       \sum_{v\in H\cap h^G}\alpha(v).
\]
For $\Gg$ the unipotent sum in $U$ equals $-1$ because $\lambda$ is nontrivial; no nonidentity semisimple belongs to $U$. Induction from $B$ can be computed by summing over fixed points: for a split element the two contributions are $\theta(a)$ and $\theta(a^{-1})$ for its eigenvalue ratio $a\ne1$; for a unipotent the unique contribution is $1$. Thus induction is unchanged under replacing $\theta$ by $\theta^{-1}$, and $P=(\pi_U-\pi_B)/2$ by induction in stages from the regular character of $B/U$. For $D$, $N'$ has $q+1$ involutions, all of sign $-1$, and a nonsplit semisimple class intersects $T'$ in two elements, both of sign $1$. For $\pi_T$ one can count ordered pairs of distinct fixed points; for $\pi_{T'}$ use the same two-element intersection. The resulting table gives each identity by substitution.
\end{proof}
\begin{center}\small
\begin{tabular}{@{}lrrrr@{}}\toprule
Character & Identity & Unipotent & Split & Nonsplit\\\midrule
$\one$ & $1$ & $1$ & $1$ & $1$\\
$\St$ & $q$ & $0$ & $1$ & $-1$\\
$P$ & $(q-2)(q+1)/2$ & $(q-2)/2$ & $-1$ & $0$\\
$D$ & $q(q-1)/2$ & $-q/2$ & $0$ & $1$\\
$\Gg$ & $q^2-1$ & $-1$ & $0$ & $0$\\
$W$ & $q^2$ & $0$ & $1$ & $1$\\
$\pi_T$ & $q(q+1)$ & $0$ & $2$ & $0$\\
$\pi_{T'}$ & $q(q-1)$ & $0$ & $0$ & $2$\\\bottomrule
\end{tabular}
\end{center}

\begin{theorem}[Optimal nonnegative character detector]\label{thm:detector}
$\Gg$ is the sum of all nontrivial complex irreducible characters of $G$, each with multiplicity one. Among genuine characters containing every nontrivial irreducible and taking only nonnegative real values, $W$ has the least possible degree, $q^2$, and is the unique minimizer.
\end{theorem}
\begin{proof}
The group $G$ is simple for $q\ge4$ \cite[Appendix]{dickson}. For any nontrivial irreducible representation, $U$ cannot act trivially: otherwise its normal closure, $G$, acts trivially. The restriction to $U$ therefore contains a nontrivial additive character. Conjugation by $T$ acts transitively on nontrivial additive characters (the nonzero scalings act freely, since multiplication by $a-1$ is bijective for $a\ne1$). Hence the chosen $\lambda$ occurs. By Frobenius reciprocity, all $q$ nontrivial irreducibles occur in $\Gg$, while the trivial one does not. Its squared norm is exactly
\[
 \frac{(q^2-1)^2+(q^2-1)}{q(q^2-1)}=q.
\]
The sum of squares of $q$ positive integer multiplicities can equal $q$ only if each is one. The least degree of any character containing the nontrivial irreducibles is thus $\Gg(1)=q^2-1$. Equality forces $\Gg$, whose value at $u$ is $-1$. Adding the trivial character gives $W$ with all values nonnegative. At degree $q^2$ the extra dimension is one, and the simple nonabelian group $G$ has no other one-dimensional character. This proves both claims.
\end{proof}
This argument avoids needing the full tabulated irreducible character table of $\PSL_2(q)$. It also shows that $P,D,\St$ partition the nontrivial irreducibles into disjoint groups, with multiplicity one. From \eqref{eq:reg} their irreducible degrees are, respectively, $q+1,q-1,q$. In particular $D$ is nonzero and contains $q/2$ irreducibles, and $P$ contains $(q-2)/2$.

\section{Transfer to Dedekind zeta functions without Artin holomorphy}
Let $M/\Q$ be any Galois extension with group $G$ as above, and set $F_H=M^H$. We assume RH for $\zeta_{\Q}$ as supplied by the upstream theorem. We use the classical meromorphic Artin $L$-function, induction and multiplicativity, and the holomorphy of nontrivial finite-order Hecke $L$-functions \cite{neukirch}. Artin induction gives $L(s,\pi_H)=\zeta_{F_H}(s)$, including the ramified Euler factors. The symbol $L(s,\lambda;F_U)$ denotes the finite-order Hecke $L$-function associated by class field theory to $\lambda$ on the abelian extension $M/F_U$, not a Dirichlet $L$-function over $\Q$.

\begin{theorem}[Exact RH transfers]\label{thm:transfer}
The following are equivalent:
\begin{enumerate}
\item RH for $\zeta_M$;
\item RH for $\zeta_{F_T}$, a zeta function of degree $q(q+1)$;
\item RH for both $\zeta_{F_B}$ and $\zeta_{F_{T'}}$, of degrees $q+1$ and $q(q-1)$;
\item $L(s,W)$ has no zeros for $1/2<\Rea s<1$;
\item the Hecke $L$-function $L(s,\lambda;F_U)$ has no zeros for $1/2<\Rea s<1$.
\end{enumerate}
More precisely, at any $\rho$ with $1/2<\Rea\rho<1$ the orders satisfy
\begin{align}
 \frac q2\ord_\rho\zeta_{F_T}&\le\ord_\rho\zeta_M
 \le(q+1)\ord_\rho\zeta_{F_T},\label{eq:ineq1}\\
 (q-1)\ord_\rho L(W)&\le\ord_\rho\zeta_M
 \le(q+1)\ord_\rho L(W).\label{eq:ineq2}
\end{align}
\end{theorem}
\begin{proof}
Let $s_0=\ord_\rho L(\St)$, $p_0=\ord_\rho L(P)$, and $d_0=\ord_\rho L(D)$. (Orders are integers, positive for zeros, negative for poles.) At this point $\zeta_\Q$ is nonzero and finite. Thus $s_0=\ord_\rho\zeta_{F_B}\ge0$. The character $P$ is a sum of induced nontrivial one-dimensional characters; its $L$-function is a product of entire Hecke $L$-functions. Similarly $D$ is induced from the nontrivial sign character $\eta$. Consequently $p_0,d_0\ge0$. No conclusion about individual Artin factors is used. Equations \eqref{eq:gamma}--\eqref{eq:reg} give
\begin{align*}
 \ord_\rho\zeta_M&=qs_0+(q+1)p_0+(q-1)d_0,\\
 \ord_\rho\zeta_{F_T}&=2s_0+p_0+d_0,\\
 \ord_\rho\zeta_{F_{T'}}&=p_0+d_0,\qquad
 \ord_\rho\zeta_{F_B}=s_0,\\
 \ord_\rho L(W)&=s_0+p_0+d_0=\ord_\rho L(s,\lambda;F_U).
\end{align*}
All the asserted equivalences and inequalities follow from positivity of the three integers and $q\ge4$. For the Dedekind zeta functions the functional equation reflects the critical strip about $1/2$; their nonvanishing on $\Rea s=1$ and in $\Rea s>1$ is classical. This converts right-half zero exclusion into RH.
\end{proof}
In particular the identities are exact meromorphic identities, not assertions that a quotient of arbitrary entire functions must be entire:
\[
 \boxed{L(s,W)=\zeta_\Q(s)L(s,\lambda;F_U)
       =\frac{\zeta_{F_B}(s)\zeta_{F_{T'}}(s)}{\zeta_\Q(s)}.}
\]
The first expression proves holomorphy apart from a simple pole at $1$; the second explains the reduction to subfields. Nonvanishing of a nontrivial finite-order Hecke $L$-function at $1$ ensures this pole is not canceled.

\begin{corollary}[Artin zero and pole exclusion off the line]\label{cor:artin}
If any of the equivalent conditions in Theorem~\ref{thm:transfer} holds, every irreducible Artin $L$-function for $M/\Q$ is holomorphic and nonzero in $1/2<\Rea s<1$. This does not assert holomorphy at zeros on the critical line.
\end{corollary}
\begin{proof}
At such a point $\rho$ let $n_\chi=\ord_\rho L(\chi)$ and form the virtual Heilbronn character $\Theta=\sum_{\chi}n_\chi\chi$. Its degree equals $\ord_\rho\zeta_M=0$. For every cyclic subgroup $C\le G$ and each one-dimensional character $\alpha$ of $C$, Frobenius reciprocity and Artin induction identify $\langle\Theta|_C,\alpha\rangle$ with $\ord_\rho L(s,\alpha;F_C)$, a nonnegative integer by Hecke holomorphy in the strip (also for $\alpha=1$, a Dedekind zeta function). Thus $\Theta|_C$ is an actual character of degree zero and vanishes. Taking $C=\langle h\rangle$ for each $h$ gives $\Theta=0$ and hence every $n_\chi=0$.
\end{proof}

\section{Optimal degree thresholds for subfield detection}
The preceding bounds are sharp in the following precisely defined sense. A family of subfields $F_{H_i}$ is \emph{spectrally detecting} if the sum of its permutation characters contains every nontrivial irreducible. This condition only tests which Artin components are visible. It does not assert that arbitrary visible meromorphic Artin factors cannot cancel. Theorem~\ref{thm:transfer} establishes the necessary analytic noncancellation for the particular optimal families, independently of Artin holomorphy.

\begin{theorem}[Sharp degrees]\label{thm:optimal}
For every $q=2^r\ge4$:
\begin{enumerate}
\item The least degree of a single spectrally detecting subfield is $q(q+1)$.
\item The least possible maximum degree of a spectrally detecting family of subfields is $q(q-1)$. Two subfields, of degrees $q+1$ and $q(q-1)$, attain it.
\end{enumerate}
These statements concern subfields of the given $M$, not auxiliary extensions or arbitrary analytic identities.
\end{theorem}
\begin{proof}
The upper bounds follow from \eqref{eq:perms} and the multiplicity-one decomposition in Theorem~\ref{thm:detector}. We prove the lower bounds by showing that any subgroup $H$ of order $>q-1$ misses at least one nontrivial character and any subgroup of order $>q+1$ misses all of $D$. Here ``misses $D$'' means $\langle D,\pi_H\rangle=\dim D^H=0$, where $D^H$ denotes the invariant space of a representation with character $D$. Since $D$ is genuine, every constituent is then missing. Similarly for $\St$.

Use Dickson's subgroup classification in the form of \cite[Theorem A.1, p.~1354]{dickson}. It lists the isomorphism types of subgroups; the argument below only uses this list and the element types of their embeddings, so no unstated conjugacy assertion is needed.

\emph{Subgroups of Borel type.} If a subgroup is isomorphic to a subgroup of $B$ with a nontrivial normal elementary abelian $2$-subgroup $V$, that subgroup fixes a projective point: every involution has a unique fixed point, commuting involutions share that point, and normality implies that the whole subgroup preserves it. Conjugate into $B$. If $V=H\cap U$ is trivial, $H$ is cyclic of order dividing $q-1$. Otherwise let the image of $H$ in $B/U$ have order $d$. Multiplication by these scalars preserves $V$. If $V\ne U$, then $V$ is a proper vector subspace over the subfield $\F_{q_0}$ generated by the image scalars (including $\F_2$ if $d=1$). Hence $|V|\le q/q_0$ and $d\le q_0-1$. Thus
\[
 |H|=|V|d\le q-q/q_0\le q-1.
\]
If $V=U$, $D^H=0$ follows from $D^U=0$: the sum of $D$ on $U$ is $q(q-1)/2-(q-1)q/2=0$.

\emph{Cyclic and dihedral types.} A cyclic subgroup of odd order has order dividing $q-1$ or $q+1$. A dihedral subgroup has order $2d$ with $d\mid q\pm1$. A proper divisor $d$ is at most $(q\pm1)/3$; for $q\ge8$ this gives $2d\le q-1$, and the same conclusion is direct for $q=4$. The only cases larger than $q-1$ are therefore the full nonsplit torus, and the full split or nonsplit dihedral groups. On a nonsplit torus, the sum of $\St$ is $q-q=0$ (all its nonidentity elements are nonsplit). For the two full dihedral groups the sums of $D$ are respectively
\[
 \frac{q(q-1)}2-(q-1)\frac q2=0,
 \qquad \frac{q(q-1)}2-(q+1)\frac q2+q=0.
\]
All these computations depend only on element orders: nontrivial odd orders dividing $q-1$ and $q+1$ are mutually exclusive. The cyclic nonsplit torus has order exactly $q+1$, so it does not affect the second lower bound.

\emph{Subfield and exceptional types.} A subgroup isomorphic to $\PSL_2(q_0)$ or $\PGL_2(q_0)$ for $q=q_0^k$, $k\ge2$, has order $q_0(q_0^2-1)$. If $k\ge3$ this is $\le q-q_0<q-1$. If $k=2$, there are $q_0^2-1=q-1$ involutions; every other nonidentity element has odd order dividing $q_0-1$ or $q_0+1$, hence dividing $q-1$. They are split in $G$. Again the sum of $D$ is zero. The exceptional $A_4$ has normal Klein four group, so lies in a Borel by the fixed-point argument. The group $S_4$ is impossible because $G$ has no elements of order four. For $A_5$, its order requires $5\mid q^2-1$, hence $r$ even. If $q=4$ it is $G$; if $q=16$ it is the just-treated square-subfield type; if $q\ge64$ its order $60$ is at most $q-1$. Finally $G$ itself has no nontrivial fixed constituents. This exhausts Dickson's list.

Since $[F_H:\Q]=g/|H|$, degrees less than $q(q+1)$ correspond to $|H|>q-1$. For a family of degrees less than $q(q-1)$ every subgroup has order $>q+1$, so all miss $D$. This proves the lower bounds.
\end{proof}
\begin{remark}
The spectral obstruction is also an obstruction to a purely formal nonnegative-order argument: assign order one to an invisible irreducible and zero to all others. Its contributions to the selected permutation characters are zero, while its regular contribution is positive. This assignment is not claimed to be realized by actual Artin zeros. No theorem excluding all other analytic methods is implied.
\end{remark}

\section{The one-sided prime-sum criterion}\label{sec:criterion}
The nonnegative detector $W$ turns the analytic end of the upstream proof into an especially transparent statement. We give the argument including the ramified Euler factors and the constant term; leaving either unspecified can lead to incorrect bounds.

\subsection{The smoothed sum and its poles}
Take a representation $V_W$ of character $W$. At each prime $p$ let $I_p$ denote inertia and $g_p$ a lift of Frobenius in a decomposition group. Put
\[
 a(p^k)=(\log p)\Tr(g_p^k\mid V_W^{I_p}),\qquad
 a(n)=0\quad\text{if $n$ is not a prime power}.
\]
At an unramified prime the trace is $W(g_p^k)$. At a ramified prime the trace is the average of $W(g_p^k h)$ over $h\in I_p$, by the projection onto invariants. It does not depend on the choice of Frobenius lift. In particular $0\le a(p^k)\le q^2\log p$. Set
\begin{align*}
 \mathcal D(s)&=-\frac{L'(s,W)}{L(s,W)}=\sum_{n\ge1}a(n)n^{-s}\qquad(\Rea s>1),\\
 c&=\FP_{s=1}\mathcal D(s)-1,\\
 S(x)&=\sum_{n\le x}a(n)(1/n-1/x),\qquad R(x)=S(x)-\log x-c.
\end{align*}
Here $\FP$ denotes the constant Laurent coefficient after subtracting $1/(s-1)$. That principal part follows from the simple pole of $L(s,W)$ at $1$. The constant $c$ is real.

\begin{lemma}[Exact Mellin transform]\label{lem:mellin}
For $\Rea s>1$,
\begin{equation}\label{eq:mellin}
 \int_1^\infty R(x)x^{-s}\,dx
 =\frac{\mathcal D(s)}{s(s-1)}-\frac{1}{(s-1)^2}-\frac{c}{s-1}.
\end{equation}
The right side continues meromorphically to $\Rea s>0$, is regular at $1$, and at any zero $\rho\ne1$ of $L(s,W)$ in that half-plane it has residue
\begin{equation}\label{eq:res}
 \frac{\ord_\rho L(W)}{\rho(1-\rho)}.
\end{equation}
There are no other poles there. At real zeros in $(0,1)$ the residue is strictly positive.
\end{lemma}
\begin{proof}
Absolute convergence for $\Rea s>1$ permits termwise integration. For any $n\ge1$,
\[
 \int_n^\infty (1/n-1/x)x^{-s}\,dx=\frac{n^{-s}}{s(s-1)}.
\]
The logarithm and constant have transforms $1/(s-1)^2$ and $c/(s-1)$. Write $\mathcal D(s)=1/(s-1)+b_0+O(s-1)$; division by $s(s-1)$ gives $1/(s-1)^2+(b_0-1)/(s-1)+O(1)$. Thus $c=b_0-1$ removes the singularity at $1$. The meromorphic continuation and the signs in \eqref{eq:res} follow because the only pole of $L(s,W)$ is at $1$, and a zero of multiplicity $m$ gives residue $-m$ in $\mathcal D$.
\end{proof}

\subsection{One-sided bounds suffice}
\begin{theorem}[One-sided RH equivalence]\label{thm:onesided}
Under RH for $\zeta_\Q$, RH for $\zeta_M$ is equivalent to the assertion that for every $\eps>0$ there exist $C_\eps>0$ and $x_\eps$ such that
\begin{equation}\label{eq:one}
 R(x)\le C_\eps x^{-1/2+\eps}\qquad(x\ge x_\eps).
\end{equation}
In the RH direction one has the stronger two-sided estimate $R(x)=O_M(x^{-1/2}\log^2 x)$ as $x\to\infty$.
\end{theorem}
\begin{proof}
Assume \eqref{eq:one}. Fix $1/2<\beta<1$ and choose $\eps=\beta-1/2$. For sufficiently large $X$ the function
\[
 H(x)=C_\eps x^{\beta-1}-R(x),\quad x\ge X,
\]
is nonnegative. Its Mellin integral converges for $\Rea s>1$, since the weights are bounded by $q^2\Lambda(n)$ and $R(x)$ grows at most as a power of $\log x$. Its meromorphic continuation to $\Rea s>\beta$ equals the negative of \eqref{eq:mellin} plus $C_\eps/(s-\beta)$ and an entire correction from the finite interval $[1,X]$.

We recall the elementary Landau principle: the Mellin transform of a nonnegative locally integrable function on $[X,\infty)$, $X>1$, has a singularity at its finite real abscissa of convergence. Indeed, if it were analytic there, take a real point $s_0$ just to the right of the abscissa so the Taylor disk extends to its left. The coefficients of the expansion in $s_0-s$ are the integrals of $H(x)x^{-s_0}(\log x)^j/j!$, all nonnegative. For positive real $s_0-s$ within the disk, monotone convergence then identifies the convergent Taylor sum with the integral to the left of the abscissa, a contradiction. (Choose $s_0$ sufficiently close using the assumed analytic neighborhood and the existing right half-plane.)

If the abscissa of $H$ were $\sigma_c>\beta$, this principle would force a real singularity there. There are no singularities at real $s\ge1$ and the only possible real poles in $(\beta,1)$ have negative residue by Lemma~\ref{lem:mellin}. A negative-residue simple pole would force the integral to tend to $-\infty$ as $s\downarrow\sigma_c$, contradicting its nonnegativity. Thus the abscissa is at most $\beta$ (or is $-\infty$). The integral is holomorphic on $\Rea s>\beta$, so \eqref{eq:mellin} has no poles there. In view of its nonzero residue at every zero of $L(s,W)$, this excludes every such zero in that region. Taking $\beta\downarrow1/2$ and applying Theorem~\ref{thm:transfer} proves RH for $\zeta_M$.

Conversely assume RH for $\zeta_M$. By Theorem~\ref{thm:transfer} and the identity for $W$, the standard prime ideal theorem under Dedekind RH \cite{grenie} for the two subfields and for $\Q$ gives, for fixed $M$,
\[
 A(t):=\sum_{n\le t}a(n)=t+E(t),\qquad E(t)=O_M(\sqrt t\log^2(2t)).
\]
Indeed logarithmic derivatives add and subtract exactly, giving the corresponding combination of the three Chebyshev functions. By partial summation $S(x)=\int_1^x A(t)t^{-2}\,dt$. Also
\[
 \mathcal D(s)=\frac{s}{s-1}+s\int_1^\infty E(t)t^{-s-1}\,dt
\]
initially for $\Rea s>1$. The integral converges near $s=1$, giving $c=\int_1^\infty E(t)t^{-2}\,dt$. Hence $R(x)=-\int_x^\infty E(t)t^{-2}\,dt=O_M(x^{-1/2}\log^2(2x))$, as claimed.
\end{proof}
The one-sided direction uses no zero-density estimate, no special exclusion of real zeros, and no auxiliary twisting hypothesis. The reason is the positive residue at a real zero, which results from using a genuine detector with known holomorphy. This is exactly the aspect of the original analytic method that transfers cleanly.

\subsection{A criterion involving rational primes only}
Take a monic integral minimal polynomial $F(X)$ of a primitive integral generator of $F_B$. Its degree is $q+1$, and its splitting field is $M$ because $B$ has trivial core. At primes not dividing its discriminant the number $f_p$ of distinct roots of $F$ modulo $p$ equals $f(g_p)$, by the usual Frobenius factorization theorem. Define $w_p$ by
\begin{equation}\label{eq:primeweights}
\begin{array}{c|cccc}
 f_p&q+1&1&2&0\\\hline
 w_p&q^2&0&1&1
\end{array}
\end{equation}
and sum only over primes $p$ not dividing the discriminant. These four possibilities exhaust the good primes. This gives a particularly simple target that requires no explicit character or representation calculation once the polynomial is specified.

\begin{corollary}[Root-count formulation]\label{cor:root}
For the above polynomial define
\begin{equation}\label{eq:rootbound}
 R_F(x)=\sum_{\substack{p\le x\\p\nmid\operatorname{disc}F}}
  w_p\log p\,(1/p-1/x)-\log x-c_F,
\end{equation}
where $c_F=c-\sum_{n\,\mathrm{omitted}}a(n)/n$ and omitted terms comprise all powers $p^k$, $k\ge2$, and first powers of the finitely many excluded primes. The defining series for the adjustment converges absolutely. Then RH for $\zeta_M$ is equivalent to $R_F(x)\le C_\eps x^{-1/2+\eps}$ eventually for every $\eps>0$.
\end{corollary}
\begin{proof}
Let $A_{\rm om}(t)$ be the sum of omitted $a(n)$ up to $t$. Bounding the number of powers $p^k\le t$, $k\ge2$, by $\sqrt t\log_2 t$, using $a(n)\le q^2\log t$, and absorbing finitely many first powers gives $A_{\rm om}(t)=O_M(\sqrt t\log^2(2t))$. Partial summation gives
\[
 \sum_{\substack{n\le x\\n\,\mathrm{omitted}}}a(n)(1/n-1/x)
 =\sum_{n\,\mathrm{omitted}}a(n)/n-\int_x^\infty A_{\rm om}(t)t^{-2}\,dt.
\]
The last integral is $O_M(x^{-1/2}\log^2(2x))$. Therefore $R_F(x)=R(x)+O_M(x^{-1/2}\log^2(2x))$; Theorem~\ref{thm:onesided} applies in both directions.
\end{proof}
One may also subtract the ordinary prime-sum remainder (already controlled by RH for $\Q$). The resulting signed contrast has weight $q^2-1$ for complete splitting, $-1$ for exactly one root, and zero for zero or two roots. In this version the logarithmic main term vanishes; the constant must still be subtracted. The weight is exactly $\Gg$ and has mean zero. The positive formulation \eqref{eq:rootbound} is preferable for geometric constructions based on divisibility. Both formulations are equivalent after the known ordinary RH bound and prime-power corrections.

\subsection{Use of full Dirichlet GRH: abelian base change}
\begin{corollary}[Transfer over abelian base fields]
Let $K/\Q$ be a finite abelian extension and $M/\Q$ as above. Under Dirichlet GRH, all the RH equivalences and sharp relative degree thresholds hold for the extension $MK/K$. The one-sided criterion holds with rational primes replaced by prime ideals of $K$, their absolute norms replacing $p$, and root counts taken in their residue fields.
\end{corollary}
\begin{proof}
The intersection $M\cap K$ is trivial by simplicity and noncommutativity, so $\Gal(MK/K)\simeq G$. The Dedekind zeta function of $K$ is a product of Dirichlet $L$-functions by the abelian factorization theorem, hence satisfies RH by the upstream result. Replace $\Q$ by $K$ throughout Theorems~\ref{thm:transfer} and \ref{thm:onesided}; the character identities and subgroup proof are unchanged, and the same nonnegative orders apply away from the line. The logarithmic derivative now sums over prime ideals $\mathfrak p$ and powers, with coefficient $\log N\mathfrak p$ times the invariant trace. Its principal part is still $1/(s-1)$, and the Mellin calculation depends only on the positive real norm. For each rational prime there are at most $[K:\Q]$ prime ideals, so the omitted prime-power estimate remains $O_{M,K}(\sqrt t\log^2(2t))$. A monic integral minimal polynomial of a generator of $(MK)^B/K$ provides the root-count criterion; $(MK)^B=M^BK$. Exclude primes dividing its discriminant ideal. This proves the asserted generalization.
\end{proof}

\section{Examples and the obstacle to a direct port}
\begin{center}
\begin{tabular}{@{}rrrrrr@{}}\toprule
$q$ & $[M:\Q]$ & Single degree & Pair degrees & $\dim W$ & $\dim\Gg$\\\midrule
4 & 60 & 20 & 5, 12 & 16 & 15\\
8 & 504 & 72 & 9, 56 & 64 & 63\\
16 & 4080 & 272 & 17, 240 & 256 & 255\\
32 & 32736 & 1056 & 33, 992 & 1024 & 1023\\\bottomrule
\end{tabular}
\end{center}
In the $q=8$ case the problem is concrete: any degree-nine polynomial with splitting group $\PSL_2(8)$ in the natural action (transitive group label 9T27) is eligible. Such fields are catalogued in the Kl\"uners--Malle number field database \cite{lmfdb}. The identity, unipotent, split and nonsplit types have respectively $1,63,216,224$ elements. Thus the weights $64,0,1,1$ have mean exactly one. This example reduces the full degree-504 problem to a single root-count inequality at rational primes, without computing the degree-504 field. We do not select or certify a particular defining polynomial or numerically bound its prime remainder here.

\begin{proposition}[Quantitative obstruction to fixed congruences]\label{prop:obstruction}
For any fixed modulus $m\ge1$, the least limiting mean square error of predicting $w_p$ by a real function of $p\bmod m$, among primes away from ramification and $m$, is exactly $q$. The unique optimal prediction on reduced classes is the constant $1$.
\end{proposition}
\begin{proof}
The simplicity and noncommutativity of $G$ imply that $M$ and the abelian cyclotomic field $\Q(\zeta_m)$ are linearly disjoint: the Galois group of their intersection is a quotient of $G$ and is abelian, hence trivial. The Galois group of the compositum is the direct product. Chebotarev's theorem \cite{neukirch} therefore gives independent uniform distributions of the conjugacy classes in $G$ (weighted by their sizes) and reduced residue classes. The average of $W$ is $1$, and $\|W-\one\|_G^2=\|\Gg\|_G^2=q$. For a prediction function $h$ on the reduced classes its limiting mean square error is
\[
 q+\frac{1}{\varphi(m)}\sum_{a\in(\mathbb Z/m\mathbb Z)^\times}(h(a)-1)^2.
\]
The second term vanishes exactly for the constant predictor. This proves the statement.
\end{proof}
In particular, even an auxiliary fixed quadratic character times a cyclotomic class function cannot exactly encode the desired nonabelian weights. Earlier reports already identify this general obstruction; the proposition supplies the sharp explicit variance for the present detector. It does not rule out genuinely nonabelian geometric counting, estimates whose modulus varies with the scale, or more elaborate analytic arguments.

\section{How to use the upstream proof method next}
The following is a proposed research program, not a completed transfer of the arithmetic argument.

\paragraph{A new arithmetic interface.} The ultimate arithmetic obligation is \eqref{eq:rootbound}. The existing Lean proof uses exact congruence counts and CM divisibilities to bound a weighted sum. To port it, one needs an arithmetic construction whose local divisibility at $p$ is governed by $W(g_p)$, or by the signed contrast after accounting for the ordinary prime sum. The finite group identities show how little nonabelian information is needed: complete splitting and the presence of exactly one fixed point. They also give the alternative expression $W=\pi_B+\pi_{T'}-1$. This could suggest using counting of primes of residue degree one in two intermediate fields, instead of building a construction directly over $M$. But the required upper bound is for their combined remainder with the correct main terms and signs, not for an arbitrary count of primes in those fields.

\paragraph{Covolume comparison.} Suppose one can realize the detector in a lattice/automorphic construction with the same three essential ingredients as the source: enough independent short vectors, a matching lower bound for integral covolumes, and an explicit error estimate. At the level of real inequalities the aim is to obtain an estimate of the form
\[
 m R_F(x)\le O_{M,\eps}(x^\eps),\qquad m\asymp \sqrt{x},
\]
after main terms, index corrections and normalizing factors cancel. This is a specification of the desired normalized conclusion, not an assertion that such lattices or inequalities have been constructed. Section~\ref{sec:criterion} proves it would be sufficient, and identifies the correct constant $c_F$. The challenge is to establish the analogues of the source's short-vector, independence, and integrality arguments with nonabelian Frobenius data. Simply substituting $w_p$ in a Lean definition does not establish them.

\paragraph{Why the smaller degrees matter.} The single-subfield route uses $\zeta_{F_T}$, whereas the pair route uses $\zeta_{F_B}$ and $\zeta_{F_{T'}}$. Any new method exploiting fields of restricted degree can focus on one of these instead of the original Galois field. Theorem~\ref{thm:optimal} proves that, within subfield-permutation spectral detection, one cannot do better than the indicated thresholds. The detector route has still smaller analytic degree, $q^2$, and its holomorphy follows from Hecke induction, not from a conjectural full Artin factorization into entire functions.

\paragraph{Further research questions.}
\begin{enumerate}
\item Can the short-form and integrality mechanisms be generalized from cyclotomic Frobenius classes to the two-point and nonsplit-torus actions of $\PSL_2(q)$? Start with $q=8$ and an explicitly certified 9T27 polynomial.
\item Is there an arithmetic construction for the signed contrast $\Gg$ that isolates its two nonzero Frobenius types, while still yielding a one-sided estimate with the required power saving? Does the field identity provide a practical way to pair local valuations?
\item For odd $q$, what are the least dimensions of nonnegative genuine characters containing every relevant irreducible, and the sharp degrees of detecting subfields? The even-characteristic argument uses transitivity on all nontrivial additive characters, which fails for $\PSL_2(q)$ in odd characteristic.
\item Which other finite simple groups admit a small nonnegative detector whose $L$-function is known to be holomorphic away from $1$ by induction from linear characters? Separating the group optimization from the analytic holomorphy requirement is essential.
\item Formalize the character identities, nonnegative-order implications, and detector optimality in Lean. The existing theorem gives the $\zeta_\Q$ hypothesis; the subgroup classification, Artin/Hecke induction, and Mellin analytic arguments would need additional formal libraries or precisely stated interfaces.
\end{enumerate}

\section{Reproducibility and limits}
The accompanying \texttt{check.py} uses only Python's standard library. It explicitly constructs $\SL_2(2^r)$ as permutations of the projective line for $r=2,3,4$, computes conjugacy classes and induction sums, and checks the displayed character identities using integer or rational arithmetic. For $q=4,8$, it also enumerates all subgroups up to conjugacy by starting at the trivial subgroup and adjoining each possible generator, using conjugacy and coset equivalence to prune only redundant extensions. The union of their conjugacy orbits consists of $59$ and $386$ subgroups respectively. It verifies the degree-obstruction conditions for every subgroup representative in both cases. For $q=16$ it checks the character identities, not the entire subgroup lattice. These finite tests corroborate the proofs, but do not establish them for arbitrary $r$; the general proofs above do.

The source manifest records the GitHub revision and the retrieved blob identifiers. No source files or previous ZIPs are redistributed in this archive. The report relies on classical Artin/Hecke and Chebotarev theorems and on Dickson's subgroup classification. It is a human-readable mathematical argument, not a Lean formalization or an independent Lean kernel audit. The literature review is targeted, not a comprehensive determination of novelty or priority. Most importantly, no proof of the arithmetic inequality \eqref{eq:rootbound} has been obtained: the additional RH and Artin zero-exclusion conclusions remain conditional on that precisely stated new input. The proven contributions are the detector construction, sharp reduction thresholds, and the exact analytic and root-count equivalences.

\begin{thebibliography}{9}
\bibitem{lean} L.~Sernau and contributors, \emph{grh-lean}, readable and compact Lean proof, revision \texttt{72fead102f663d5c7a407964a13dc0cf9d977207}. \url{https://github.com/lsernau/grh-lean/tree/72fead102f663d5c7a407964a13dc0cf9d977207}.
\bibitem{incoming} V.~Reshetnikov, \emph{ProveIt}, \texttt{docs/incoming}, GRH transfer archives, consulted November 25, 2026; see accompanying source manifest. \url{https://github.com/VladimirReshetnikov/ProveIt/tree/main/docs/incoming}.
\bibitem{dickson} R.~M.~Guralnick and M.~E.~Zieve, Polynomials with PSL(2) monodromy, \emph{Annals of Mathematics} \textbf{172} (2010), 1315--1359, especially Appendix A, Theorem A.1 (Dickson's classification), p.~1354. \url{https://annals.math.princeton.edu/wp-content/uploads/annals-v172-n2-p11-p.pdf}.
\bibitem{neukirch} J.~Neukirch, \emph{Algebraic Number Theory}, Springer, 1999, Chapter VII (zeta and $L$-functions, Artin induction, prime distribution); classical class field theory for the one-dimensional factors.
\bibitem{grenie} L.~Greni\'e and G.~Molteni, Explicit versions of the prime ideal theorem for Dedekind zeta functions under GRH, \emph{Mathematics of Computation} \textbf{85} (2016), 889--906. \url{https://doi.org/10.1090/mcom3031}. Only the standard fixed-field GRH Chebyshev error estimate is used here.
\bibitem{lmfdb} J.~Kl\"uners and G.~Malle, \emph{A Database for Number Fields}, transitive groups of degree 9, entry 9T27. \url{https://galoisdb.math.upb.de/groups?deg=9&sort=center}.
\end{thebibliography}
\end{document}
